%REVIEW-ADDR: W1-zvf-tautology
%REVIEW-ADDR: Q1-zvf-formalization
\section{Formal Definition of the Zero-Variance Fraction, Non-Tautology, and Scope}
\label{sec:appendix:zvf-formalization}

\subsection{Formal Definition}

Let a GRPO batch $B_t$ at optimizer step $t$ consist of $|B_t|$ prompt-groups,
each group $g$ containing $G$ rollouts with scalar outcome rewards
$r_{g,1}, \ldots, r_{g,G} \in \mathbb{R}$. We define the \emph{Zero-Variance
Fraction} (ZVF) as
\begin{equation}
\label{eq:zvf}
\mathrm{ZVF}_t \;=\; \frac{1}{|B_t|} \sum_{g \in B_t}
  \mathbb{1}\!\left[\;\widehat{\mathrm{Var}}(r_{g,1}, \ldots, r_{g,G}) = 0\;\right],
\end{equation}
where $\widehat{\mathrm{Var}}$ is the unbiased sample variance and $G\ge2$.
This exact statistic counts identical rewards and implies zero centered
reward coefficients, not necessarily zero total-objective gradients.
A separate near-constant telemetry statistic, used only for logging and
never as a gradient mask, uses variance $\le10^{-6}$;
for real rewards that threshold does not imply zero coefficients. For example,
$(0,0,0,10^{-4})$ has sample variance $2.5\times10^{-9}$ but nonzero centered
rewards. For binary rewards the smallest mixed-group sample variance is
$1/G$; hence the tolerance agrees with exact collapse only when $1/G>10^{-6}$,
as for the reported group sizes. The expectation below concerns exact collapse.

\subsection{Why ZVF Is Not a Tautological Re-expression of Mean Reward}

A reviewer may reasonably suspect that under binary outcome rewards
$r \in \{0,1\}$ and group-relative advantages, $\mathrm{ZVF}_t$ is a direct
function of the batch mean reward $\bar{r}_t$: if every group succeeds
(or fails) uniformly, $\mathrm{ZVF}_t = 1$. This is true at the extremes
$\bar{r}_t \in \{0,1\}$ but does \emph{not} hold at intermediate values.
For $G$ i.i.d.~Bernoulli rollouts with prompt-specific success probability
$p_g$ and prompts $g \sim \mathcal{D}$, the expected ZVF under a null independence model is
\begin{equation}
\mathbb{E}[\mathrm{ZVF}_t] \;=\; \mathbb{E}_{g \sim \mathcal{D}}\!\left[\,p_g^{G} + (1-p_g)^{G}\,\right].
\end{equation}
The right-hand side depends on the full prompt-difficulty distribution
$\{p_g\}$, not only on its first moment. Two populations can therefore
share the same mean reward $\bar r = 0.5$ yet have very different ZVF:
a bimodal population with $p_g \in \{0,1\}$ yields $\mathrm{ZVF}=1$,
while a uniform-hard population with all $p_g = 0.5$ and $G=8$ yields
$\mathrm{ZVF} \approx 2 \cdot (0.5)^8 \approx 0.008$. The substantive
claim is thus not ``high mean reward implies high ZVF,'' but rather that
ZVF is sensitive to \emph{how} success mass is distributed across prompts.

\subsection{What the Released Artifact Establishes}

The released artifact bundled with this paper contains aggregate reward
trajectories, held-out checkpoint summaries, and the cross-framework parser
specification released with the anonymous artifact. It does
\emph{not} contain the per-group step-level rollout logs needed to support
a new partial-correlation table that controls simultaneously for
batch-mean reward, entropy, KL drift, and advantage variance. We therefore
do \emph{not} claim a fresh incremental-$R^2$ estimate here, and we mark
any implication that the current artifact proves a new partial-correlation
result beyond the already released summary traces as rerun required.

The narrower empirical claim made in the revised paper is the one directly
supported by the released evidence: in outcome-reward GRPO, high ZVF
coincides with the disappearance of within-group contrast, and at any step
where $\mathrm{ZVF}=1$ the group-relative advantage is mechanically zero, so
the centered reward-contrast term is zero in the stated objective. We additionally observe co-occurrence
between elevated ZVF and reward plateaus or regressions in some runs, but
we do not establish a causal link: the within-task stratified analysis in
Section~\ref{sec:zvf-analysis} does not support a within-task predictive
relationship between mean ZVF and last-10 reward. The operational meaning
of ZVF in the present benchmark is therefore the local mechanical claim,
not a performance-causal one.

\subsection{Boundary Conditions and Scope}

We state explicitly when ZVF is \emph{not} informative, partly in response
to the tautology concern:
\begin{itemize}[leftmargin=1.2em,itemsep=0pt]
  \item \textbf{Dense or continuous rewards.} When rewards are continuous
        (e.g.,~process-reward models or graded code-execution partial credit),
        $\widehat{\mathrm{Var}}$ is almost surely nonzero and
        $\mathrm{ZVF} \to 0$. ZVF degenerates and must be replaced by a
        per-step-reward-variance diagnostic or alternative diagnostic surrogates.
  \item \textbf{$G=1$.} The unbiased sample variance and this statistic are undefined; exclude such groups.
  \item \textbf{SFT-saturated baselines.} When the policy already solves
        the task ($\bar p_g \to 1$ for essentially every prompt),
        $\mathrm{ZVF} \to 1$ and the diagnostic loses discriminative power;
        in this regime RL no longer provides useful signal anyway.
  \item \textbf{Non-outcome-reward RL.} For DPO- or regression-style
        training there are no rollout groups and ZVF is ill-defined.
\end{itemize}
Within its intended scope -- outcome-reward GRPO on verifiable tasks with
$G \ge 4$ and non-saturated reward support -- ZVF is a cheap descriptive statistic; its early-warning utility is unvalidated. Outside that scope, the paper now recommends
explicit surrogates instead of over-claiming portability.
